% 03-introduction.tex — Introduction
% Structure: Clinical background & auto-FMT -> Methodological challenges in longitudinal microbiome analysis -> Study aims & research questions.

\section{Introduction}

Allogeneic hematopoietic cell transplantation (allo-HCT) is a potentially curative therapy for hematologic malignancies and bone marrow failure syndromes, yet the conditioning chemotherapy, radiation, and broad-spectrum antibiotics administered during the peri-transplant period severely disrupt the host gut microbiome. This acute ecological injury is characterized by a steep collapse in microbial richness and the opportunistic domination of the intestinal lumen by pathobionts such as \textit{Enterococcaceae}, and this loss of gut microbiota diversity during allo-HCT is associated with increased mortality \citep{taur2018reconstitution}. Autologous fecal microbiota transplantation (auto-FMT)---in which a patient's own pre-treatment stool is cryopreserved prior to conditioning and reinfused following antibiotic cessation and hematopoietic engraftment---offers a targeted strategy to restore baseline microbial diversity without exposing immunocompromised recipients to third-party donor pathogens \citep{herman2024assessing,herman2024assessingb}. In a randomized trial of 25 allo-HCT patients (14 auto-FMT, 11 control), auto-FMT was reported to boost microbial diversity and reestablish the intestinal microbiota composition present before antibiotic treatment and transplantation \citep{taur2018reconstitution}. However, determining whether and how auto-FMT alters post-transplant microbiome trajectories requires analytical frameworks capable of separating intervention-associated recovery from the massive background temporal shifts induced by transplantation itself.

Evaluating microbiome interventions in longitudinal clinical cohorts presents several interrelated statistical and bioinformatic challenges. First, clinical microbiome trials collect repeated fecal specimens from the same individuals across irregular follow-up schedules; pooling these observations into cross-sectional rank tests violates sample independence (pseudoreplication), inflates false discovery rates, and allows large pre- versus post-transplant variance to mask between-group trajectory differences \citep{herman2024assessingb,zhou2021linda,zhang2026bayesian}. Second, standard diversity workflows commonly normalize library sizes via a single random subsampling without replacement (rarefying) to a fixed read threshold, a practice that discards valid sequencing reads and excludes low-depth specimens that often correspond to the most severely depleted post-antibiotic states \citep{mcmurdie2013waste,raspet2024facilitating}. Third, while linear mixed-effects (LME) models account for within-subject correlation through patient-specific random effects, standard restricted maximum likelihood (REML) optimizers frequently encounter convergence failures under small sample sizes and unbalanced designs \citep{barrera2025stochastic}, and single-slope linear specifications across pre- and post-intervention time points can oversimplify piecewise or time-varying recovery dynamics \citep{fontaine2026locally}. Finally, general diversity recovery or temporal taxonomic shifts identified by supervised machine learning do not by themselves establish the engraftment of reinfused baseline strains \citep{herman2024assessing,herman2024assessingb} or correct for the compositional nature of sequencing count data \citep{zhou2021linda}.

In this study, we performed a structured longitudinal reanalysis of 16S rRNA amplicon sequencing data from the randomized auto-FMT trial NCT02269150 (356 retained fecal samples from 24 allo-HCT patients across days $-10$ to $+70$ relative to transplantation) using the QIIME 2 platform. Specifically, we addressed three research questions: (1) how gut microbial richness and community structure change across the peri-transplant timeline; (2) whether longitudinal LME modeling reveals a divergence in richness trajectories between the auto-FMT and control arms that is obscured when samples are pooled in cross-sectional Kruskal-Wallis tests; and (3) which genus-level bacterial taxa systematically track post-transplant temporal succession in random forest regression. Beyond reporting these empirical findings, we use this reanalysis as a methodological case study to examine six key analytical limitations---small unbalanced sample sizes, LME non-convergence, single-slope temporal specification, single-step rarefying, repeated-measures pseudoreplication, and the distinction between temporal succession and microbiome engraftment---and map each limitation to emerging statistical solutions in the microbiome literature.
